\documentclass[10pt,a4paper]{amsart}
\usepackage[margin=19mm]{geometry}
\usepackage{amsmath,amssymb,mathtools}
\usepackage[scr=rsfs,cal=euler]{mathalfa}
\usepackage{xfrac}
\usepackage[unicode,hidelinks,bookmarks=false]{hyperref}
\newcommand{\ie}{i.e.\ }
\newcommand{\cf}{cf.\ }
\newcommand{\resp}{respectively\ }
\newcommand{\Subsection}{\subsection}
\providecommand{\nobreakdash}{\nobreak\hbox{-}}
\setlength{\emergencystretch}{2em}
\multlinegap0pt
\usepackage{amssymb}
\usepackage{mathtools}
\usepackage{xcolor}
\newtheorem{lemma}{Lemma}
\newcommand{\GG}{\mathcal{G}}
\newcommand{\Hil}{\mathsf{H}}
\title{RFD property for groupoid C*-algebras of amenable groupoids and for crossed products by amenable actions}
\author{Tatiana Shulman, Adam Skalski}
\date{Source: arXiv:2305.12122v3 (CC BY 4.0)}
\begin{document}
\maketitle

\noindent\emph{Source and licence.} RFD property for groupoid C*-algebras of amenable groupoids and for crossed products by amenable actions. Version arXiv:2305.12122v3 (CC BY 4.0), distributed by arXiv under the Creative Commons Attribution 4.0 International licence, http://creativecommons.org/licenses/by/4.0/, as recorded in the arXiv licence metadata for this exact version and shown on the abstract page https://arxiv.org/abs/2305.12122v3. Full licence notice: \texttt{licenses/arxiv-2305.12122-CC-BY-4.0.txt}.\par\smallskip

\noindent\textit{Editorial scope.} Counted unit: the article's lemma lem:concreteform (source0.tex lines 296--302). The article's complete original proof of it is delivered with it (source0.tex lines 304--339, one \texttt{proof} environment of 36 lines). No mathematics is written, repaired or re-derived: the statement and the proof are the article's own lines, in the article's own notation, and no retained line is substituted or re-flowed.  No further source-local context is needed: every cross-reference of every retained line resolves inside the two delivered spans.  Nothing was omitted from any retained span, so the package carries no omission record.  No retained line contains a citation, so the wrapper carries no bibliography and no bibliography entry.  No retained line contains a figure environment, an image-inclusion command or drawing code, so no image is delivered and the package declares no asset dependency.  Editorial formatting only, in addition: an AMS \texttt{amsart} preamble that restates the article's own declarations and macros used by the retained lines, namely the environments \texttt{lemma} on separate counters, together with the standard packages those lines need.  The retained environments print in this document as Lemma 1 in order of appearance, whereas the article prints the same items with the numbering of its own full document.  The complete native source of the article as distributed in the arXiv e-print for this exact version is bundled unchanged as \texttt{sources/141792/source0.tex}.
\begin{lemma}\label{lem:concreteform}
Suppose that $F\subset X$ is a finite orbit of the form $\{x_1, \ldots, x_n\}$, where $n = |F|$. For all $i,j=1, \ldots, n$ write $\GG_{j}^i:= \GG_{x_j}^{x_{i}}$, set $\gamma_{1}=e_{\GG_1^1}$   and for $i=2,\ldots, n$ fix  an element $\gamma_i \in \GG_{1}^{i}$, writing $\overrightarrow{\gamma}=(\gamma_1, \ldots, \gamma_n)$. Consider an arbitrary representation $\rho: \GG_{1}^1 \to B(\Hil_1)$, where $\Hil_1$ is a Hilbert space. Further assume we have a collection of unitaries ${\bf{z}}:=(z_i)_{i=1}^{n}$ such that $z_1=1_{H_i}$ and $z_i:\Hil_1\to \Hil_{i}$ for $i=2,\ldots, n$, where $\Hil_2=\cdots=\Hil_n = \Hil_1$. Finally
introduce a new Hilbert space $\mathcal{H}:=\bigoplus_{i=1}^n \Hil_i$. Then the `block-matrix' formula
\[ (\pi_{F, \rho, \bf{z}, \overrightarrow{\gamma}} (f))_{i,j} = \sum_{g \in \GG_j^i} f(g) z_{i} \rho (\gamma_{i}^{-1}g \gamma_{j}) z_j^*\in B(\mathcal{H}),  \]
$f \in C_c(\GG)$, $i,j=1, \ldots, n$,
defines a representation of $C^*(\GG)$ on $B(\mathcal{H})$.
	\end{lemma}

\begin{proof}
The proof is a simple check. Choose $(F, \rho, \bf{z}, \overrightarrow{\gamma})$ as above and  write $\pi:=\pi_{F, \rho, \bf{z}, \overrightarrow{\gamma}}$.
	
First note that each of the sums in the displayed formula is in fact finite, as $\GG$ is \'etale, so each $\GG_j^i$ is discrete (and $f$ is compactly supported). The correspondence $f\to \pi(f)$ is clearly linear.

Fix $i,j=1, \ldots, n$. We have  $(\pi(1_X))_{i,j}= \delta_{i,j} z_i \rho(\gamma_i^{-1}\gamma_i) z_i^* = \delta_{i,j} 1_{H_i}$. If $f \in C_c(\GG)$ then
\begin{align*}  (\pi (f^*))_{i,j} &= \sum_{g \in \GG_j^i} f^*(g) z_{i} \rho (\gamma_{i}^{-1}g \gamma_{j}) z_j^*
= \sum_{g \in \GG_j^i} \overline{f(g^{-1})} z_{i} \rho (\gamma_{i}^{-1}g \gamma_{j}) z_j^*
\\& = \sum_{g \in \GG_i^j} \overline{f(g)} z_{i} \rho (\gamma_{i}^{-1}g^{-1} \gamma_{j}) z_j^* =
 \sum_{g \in \GG_i^j} \overline{f(g)} z_{i} \rho ((\gamma_{j}^{-1}g \gamma_{i})^{-1}) z_j^*
\\&=
\sum_{g \in \GG_i^j} \overline{f(g)} z_{i} \rho ((\gamma_{j}^{-1}g \gamma_{i}))^* z_j^*
=
 \left( \sum_{g \in \GG_i^j} f(g) z_{j} \rho ((\gamma_{j}^{-1}g \gamma_{i})) z_i^* \right)^*
\\&= \left((\pi (f))_{j,i}\right)^*, \end{align*}
so that $\pi (f^*) = \left(\pi (f)\right)^*$.
Finally if $f, h\in C_c(\GG)$ and $g \in \GG_j^i$ then we have, setting $\GG^i=\{g \in \GG: r(g) =x_i\}$,
\[ (f\star h)(g) = \sum_{\kappa \in \GG^i} f(\kappa)h(\kappa^{-1}g) = \sum_{k=1}^n
\sum_{\kappa \in \GG_k ^i} f(\kappa)h(\kappa^{-1}g),\]
as if $r(\kappa)= x_i$ then $s(\kappa) \in F$. Thus
\begin{align*} (\pi (f\star h))_{i,j} &= \sum_{g \in G_j^i} (f \star h)(g) z_{i} \rho (\gamma_{i}^{-1}g \gamma_{j}) z_j^*
= \sum_{g \in G_j^i} \sum_{k=1}^n
\sum_{\kappa \in \GG_k^i} f(\kappa)h(\kappa^{-1}g) z_{i} \rho (\gamma_{i}^{-1}g \gamma_{j}) z_j^*.
\end{align*}
On the other hand
\begin{align*} (\pi (f)  \pi (h))_{i,j}
	&= \sum_{k=1}^n 	(\pi_{F, \rho, \bf{z}} (f))_{i,k}  (\pi_{F, \rho, \bf{z}} (h))_{k,j}
	\\&= \sum_{k=1}^n \sum_{g \in \GG_k^i} f(g) z_{i} \rho (\gamma_{i}^{-1}g \gamma_{k}) z_k^*
	\sum_{g'\in \GG_j^k} h(g') z_{k} \rho (\gamma_{k}^{-1}g' \gamma_{i}) z_j^*
	\\&=
\sum_{k=1}^n \sum_{g \in \GG_k^i} 	\sum_{g'\in G_j^k} f(g) h(g')  z_i \rho (\gamma_{i}^{-1}g \gamma_{k}) \rho (\gamma_{k}^{-1}g' \gamma_{i}) z_j^*
\\&=
\sum_{k=1}^n \sum_{g \in \GG_k^i} 	\sum_{g'\in G_j^k} f(g) h(g')  z_i \rho (\gamma_{i}^{-1}g g' \gamma_{i}) z_j^*.
\end{align*}	
A straightforward reindexing argument shows that the two  sums displayed above coincide.
\end{proof}

\end{document}
